\documentclass[11pt]{article}
\usepackage[margin=1in]{geometry}
\usepackage{amsmath,amssymb,amsthm,booktabs,microtype,needspace}
\usepackage[colorlinks=true,urlcolor=blue,citecolor=blue,linkcolor=blue]{hyperref}
\newtheorem{theorem}{Theorem}
\newtheorem{lemma}{Lemma}
\newtheorem{proposition}{Proposition}
\newcommand{\Lmin}{\operatorname{Lmin}}
\newcommand{\Lmax}{\operatorname{Lmax}}
\title{Fair core selections:\\a sharp boundary between pairs and triples}
\author{Maxwell Flanagan}
\date{5 September 2026}
\begin{document}
\maketitle
\begin{abstract}
We compare full leximin and descending-lexicographic-minimum core payoffs
in unit-capacity matching markets. For every graph matching game with nonempty
core, leximin guarantees half of each player's maximum core payoff and the
sharp normalized distance is $\lfloor n/2\rfloor/(\lfloor n/2\rfloor+1)$ for
$n\ge4$. Three-sided assignment games instead permit distance tending to two
and no positive dimension-independent individual guarantee, already with
weights in $\{0,1,2,3\}$. Complete additive three-sided markets with equal sector sizes retain
a one-third guarantee and a sharp distance supremum of $4/3$.
The proofs use the entire core and count every player. A priority audit
identifies the bipartite half-guarantee as a short consequence of existing
algorithmic invariants; the sharp distance comparison and sparse triple
construction are the principal candidate contributions. External priority
and significance are not established.
\end{abstract}

\section{One question, three market classes}
The two fairness rules are established objects: leximin raises the lowest
payoffs successively, while leximax-min lowers the highest successively.
Older egalitarian work and recent combinatorial algorithms motivate their
comparison \cite{klijn,vazirani}. We ask how much the choice of rule can alter
who receives a fixed surplus, and whether leximin protects each player's
share of their own best stable outcome.

A unit-capacity matching game has a finite player set $N$ and nonnegative
weights on feasible teams. A coalition earns the maximum weight of disjoint
teams contained in it. Teams are edges of a graph in the pairwise model, and
triples containing one player from each of three sectors in the three-sided
model. Let $T>0$ be the \emph{integral} matching value. The entire core is
\begin{equation}\label{eq:core}
 C=\{p\ge0:\ \textstyle\sum_{i\in N}p_i=T,\quad
                    \sum_{i\in e}p_i\ge w_e\text{ for every team }e\}.
\end{equation}
These inequalities imply all coalition inequalities by summing over an
optimal coalition matching. We assume $C\ne\varnothing$. Bipartite graphs
ensure this, whereas general graphs and triple arrays need not.

Write $x=\Lmin(C)$ for the lexicographic maximum of ascending sorted payoffs,
$z=\Lmax(C)$ for the lexicographic minimum of descending sorted payoffs, and
\[
 \rho(C)=\frac{\|x-z\|_1}{T},\qquad
 H_i(C)=\max_{y\in C}y_i.
\]
All players count toward $n=|N|$, including fixed-zero players. These are full
selectors, not arbitrary first-stage maximin or minimax points. If $H_i=0$,
any fractional guarantee for that coordinate is vacuous.

The quantity $\rho/2$ has a direct accounting interpretation:
\[
 \frac{\rho}{2}=\frac1T\sum_{i:z_i>x_i}(z_i-x_i).
\]
It is the minimum fraction of total surplus that must change recipients when
moving between the two allocations. Thus the two-sided bound approaches
one-half of the surplus, whereas three-way cooperation can require moving
almost all of it. This is a worst-case statement, not an empirical claim.

\begin{theorem}\label{thm:sharp}
For any unit-capacity graph matching game with nonempty core, fix an optimal
matching and let $k\ge1$ count its positive-weight edges. Then
\[
 x_i\ge H_i(C)/2\quad\hbox{for every }i,\qquad
 \|x-z\|_1/T\le k/(k+1).
\]
The distance bound holds for any minimax core point $z$. For $k=1$ the full
selectors coincide. For each $k\ge2$ a bipartite market attains the bound.
Consequently the exact $n$-player maximum is zero for $n=2,3$ and
$\lfloor n/2\rfloor/(\lfloor n/2\rfloor+1)$ for $n\ge4$, both for bipartite
markets and for all nonempty-core graph matching games.
\end{theorem}

\begin{theorem}\label{thm:main}
For all integers $L\ge1,R\ge0$, a three-sided assignment game with nonempty
entire core and weights in $\{0,1,2,3\}$ has
\[
 n=9(L+R)+12,\qquad T=3(L+R+2),
\]
\begin{equation}\label{eq:rho}
 \rho(L,R)=\frac{2((R+2)2^L-1)}{(2^L+2)(L+R+2)}.
\end{equation}
A player has $H_i=1$ and $x_i=2/(2^L+2)$. Hence no positive
 dimension-independent individual fraction survives, and $R=L^2$ gives
$\rho\to2$, the sharp supremum. The family has equally sized sectors and
only $12(L+R)+9$ positive triples.
\end{theorem}

\begin{center}
\small
\begin{tabular}{lll}
\toprule
Entire-core domain & Uniform $x_i/H_i$ guarantee & Distance supremum\\
\midrule
Graph matching, nonempty core & $1/2$ & $1$\\
Equal-sized complete additive triples & $1/3$ & $4/3$\\
Arbitrary three-sided, nonempty core & none positive & $2$\\
\bottomrule
\end{tabular}
\end{center}
The middle row is proved below. It shows why the additive benchmark does not
settle the unrestricted question. The graph extension shows that
bipartiteness itself is not the transfer mechanism. These statements concern
unit capacity and do not assert analogous results for $b$-matching games.

\section{Pairwise transfer and half domination}
Summing the inequalities of any optimal positive matching in \eqref{eq:core}
shows that each matched pair has its designated total $s_i$ and every
unmatched player is fixed at zero.
\begin{lemma}\label{lem:half}
For every $y\in C$, $y_i\le2x_i$ for every player $i$. In particular
$\|x-y\|_1/T\le1$ for every core allocation $y$.
\end{lemma}
\begin{proof}
Put $d=y-x$ and suppose $d_a=\delta>0$. Define global vertex sets
\[
 A=\{i:d_i\ge\delta\},\qquad B=\{i:d_i\le-\delta\}.
\]
Matched-pair changes are opposite and unmatched changes are zero, so matching
edges biject $A$ and $B$. At a tight edge $uv$, feasibility of $y$ gives
$d_u+d_v\ge0$. If $u\in B$ and $v\notin A$, then
$d_u+d_v<0$, a contradiction. Therefore the direction
$h=\mathbf1_A-\mathbf1_B$ preserves the total and harms no tight inequality.
Every decreasing coordinate has $x_i\ge\delta$ because $y_i\ge0$.
Finiteness of the slack constraints yields $x+\epsilon h\in C$ for all
sufficiently small $\epsilon>0$.

If $\min_Bx>\min_Ax$, this step raises the smallest affected payoffs while
lowering only strictly larger ones, contradicting full leximin even when
smaller fixed payoffs occur elsewhere. Hence
\[
 \delta\le\min_Bx\le\min_Ax\le x_a,
 \qquad y_a=x_a+\delta\le2x_a.
\]
The other coordinates follow by repeating the argument or by $y_i\le x_i$.
The overlap is at least $T/2$, giving the distance statement.
No bipartition was used.
\end{proof}
The fraction is sharp on a single edge. In the bipartite case it also follows
from Vazirani's algorithmic invariants; Section~\ref{sec:priority} makes this
priority qualification explicit. The proof above is a direct geometric proof
and also handles nonbipartite nonempty-core games.

\section{The sharp finite bound and its attainment}
\begin{proof}[Proof of Theorem~\ref{thm:sharp}: upper bound]
On matched pair $i$, write the change from $x$ to $z$ as $(d_i,-d_i)$,
and put
\[
 q_i=|d_i|,\quad D=\sum_iq_i,\quad M=\max_a z_a>0.
\]
Thus $\|x-z\|_1=2D$. Both $x$ payoffs in pair $i$ are at least $q_i$:
the decreasing endpoint by nonnegativity of $z$, and the increasing
endpoint by Lemma~\ref{lem:half}. The increasing endpoint's $z$ payoff
is at least $2q_i$. Therefore
\begin{equation}\label{eq:pair}
 s_i\ge2q_i,\qquad q_i\le M/2,\qquad D\le kM/2.
\end{equation}
These statements also hold when $q_i=0$.

Some player $a$ satisfies $z_a=M$ and $x_a\ge M$. Otherwise, a
sufficiently short move from $z$ toward $x$ lowers every maximum
coordinate and leaves every other coordinate below $M$, contradicting
minimax optimality. Since $M>0$, this player belongs to a positive
matched pair $p$. Its displacement is $q_p=x_a-M$; the other $x$
payoff in that pair is at least $q_p$. Thus $s_p\ge M+2q_p$.
Summing this strengthened inequality with \eqref{eq:pair} for the other
pairs gives
\[
 T\ge2D+M\ge2D\left(1+\frac1k\right).
\]
This proves the bound, also when $D=0$.

If $k=1$, all other payoffs are fixed at zero. The remaining two sum
to $T$ and range over a closed interval of splits. Maximizing their
minimum is the same as minimizing their maximum, and uniquely selects
the feasible split closest to equality. Hence $x=z$ in this case.
\end{proof}
\begin{proof}[Attainment and player count]
For $k\ge2$, take a complete $k\times k$ market with identical rows
\[
 (w_{i1},\ldots,w_{ik})=(1,\ldots,1,2).
\]
Every perfect matching has value $T=k+1$, and every edge belongs to one.
Complementarity consequently makes every edge inequality tight throughout
the core. Writing $c=(1,\ldots,1,2)$, the entire core is exactly
\[
 u_i=t\quad(1\le i\le k),\qquad
 v_j=c_j-t\quad(1\le j\le k),\qquad0\le t\le1.
\]
Its minimum coordinate is $\min(t,1-t)$, uniquely maximized at $t=1/2$.
Its maximum is $2-t$, uniquely minimized at $t=1$. Every coordinate
changes by $1/2$, so $\|x-z\|_1=k$ and $\rho=k/(k+1)$.
Adding one isolated player leaves the full selectors otherwise unchanged.
Every matching on $n$ players has at most $\lfloor n/2\rfloor$ edges,
which proves the exact $n$-player formula.
\end{proof}


\section{An explicit market}
All color indices are modulo three. Introduce global players $G_i$ of color
$i$. At levels $\ell=0,\ldots,L$ introduce $U_{\ell i},V_{\ell i},D_{\ell i}$
for $i=0,1,2$. At even levels their colors are $i,i+1,i+2$; at odd levels
the colors of $U$ and $V$ are interchanged. Put $s_0=2$ and $s_\ell=1$
for $\ell\ge1$. Assign the following positive triple values:
\begin{center}
\begin{tabular}{ll}
\toprule
Triple & Value\\
\midrule
$\{U_{\ell i},V_{\ell i},D_{\ell i}\}$ & $s_\ell$\\
$\{U_{\ell i},V_{\ell i},G_{i+2}\}$ & $s_\ell$\\
$\{U_{0i},G_{i+1},G_{i+2}\}$ & $1$\\
$\{U_{\ell i},V_{\ell-1,j},V_{\ell-1,k}\}$, $\ell\ge1$ & $2$\\
$\{V_{\ell i},U_{\ell-1,j},U_{\ell-1,k}\}$, $\ell\ge1$ & $3$ if $\ell=1$, else $1$\\
\bottomrule
\end{tabular}
\end{center}
In the last two rows, $\{j,k\}=\{0,1,2\}\setminus\{i\}$.
Each listed triple has exactly one player of each color.

For every $i$ introduce $R$ copies $(U'_{ri},V'_{ri},D'_{ri})$ with the same
colors as its top-level triple. Give value $1$ to each copy's local and global
pair triples $\{U'_{ri},V'_{ri},D'_{ri}\}$ and
$\{U'_{ri},V'_{ri},G_{i+2}\}$, and to its cross-triples
\[
 \{U'_{ri},V_{Li},G_{i+2}\},\qquad
 \{V'_{ri},U_{Li},G_{i+2}\}.
\]
Every unlisted triple has value zero. There are $3(L+R)+4$ players per sector.

\section{Integral value and entire core}
The triples using each pair's own $D$ form a disjoint integral matching with
weight $6+3L+3R=T$. Consider the allocation assigning $1$ to both root
$U$ and root $V$, $0$ to every subsequent $U$, $1$ to every subsequent $V$,
the analogous values to copies, and $0$ to every dummy. It has total $T$
and covers every weighted triple. Summing its payoffs over any disjoint
matching bounds that matching's weight by $T$. Hence the displayed integral
matching is optimal and this allocation belongs to the entire core.

\Needspace{12\baselineskip}
\begin{lemma}\label{lem:core}
Every core allocation, and only a core allocation, has all dummy payoffs zero
and the following form:
\begin{align*}
 U_{0i}&=1+q_{0i}, & V_{0i}&=1-q_{0i},\\
 U_{\ell i}&=q_{\ell i}, &V_{\ell i}&=1-q_{\ell i}\quad(\ell\ge1),\\
 U'_{ri}&=q_{Li}, &V'_{ri}&=1-q_{Li},
\end{align*}
where $0\le q_{\ell i}\le1$ and
\begin{equation}\label{eq:rec}
 q_{\ell i}=q_{\ell-1,j}+q_{\ell-1,k},\qquad
 q_\ell=(J-I)q_{\ell-1}.
\end{equation}
Here $J$ is the $3\times3$ all-ones matrix.
\end{lemma}
\begin{proof}
Sum the inequalities of the displayed matching. Efficiency and nonnegativity
force every unmatched $G$ to zero and every matching triple to equality.
The alternative pair triple using $G$ then forces its local $D$ to zero and
its pair to total $s_\ell$. Root lower-bound triples give $U_{0i}\ge1$.
At level one, the two recurrence inequalities become respectively
$q_{1i}\ge q_{0j}+q_{0k}$ and $q_{1i}\le q_{0j}+q_{0k}$ after substituting
the root offsets. At subsequent levels the same substitutions give these
opposite inequalities without offsets. Cross-triples give opposite inequalities
between a copied $U$ and its original top $U$, hence equality.
Nonnegativity supplies the stated bounds.

Conversely, the parameterization satisfies every displayed triple inequality,
has total $T$, and is nonnegative. All omitted zero-weight triple inequalities
are automatic. It therefore satisfies every coalition inequality by summing
over a matching. This proves equality with the entire core.
\end{proof}

\section{The full selectors and quantitative separation}
For minimax, every root $U$ is at least $1$. The allocation just used, denoted
$z$, has maximum payoff $1$. Any core allocation with maximum at most $1$
must have $U_{0i}=1$ for all $i$, hence $q_0=0$. Lemma~\ref{lem:core} then
fixes every coordinate. Thus $z$ is the unique minimax allocation and the
full descending-lexicographic minimum.

For leximin, all $D,G$ are fixed at zero. Every other coordinate can
simultaneously be positive by taking $q_0=t(1,1,1)$ with $0<t<2^{-L}$.
Consequently the full procedure first fixes exactly the dummy coordinates
at zero, then maximizes a common floor $\eta$ for all remaining coordinates.
If such a floor is feasible, level-one $U$ and top-level $V$ imply
\[
 2\sum_iq_{0i}=\sum_iq_{1i}\ge3\eta,\qquad
 2^L\sum_iq_{0i}=\sum_iq_{Li}\le3(1-\eta).
\]
Combining gives $\eta\le2/(2^L+2)$. Equality is attained by
\begin{equation}\label{eq:x}
 t=\frac1{2^L+2},\qquad q_{\ell i}=2^\ell t.
\end{equation}
The least nonroot $U$ is $2t$ and the least $V$ is $1-2^Lt=2t$;
root and copy payoffs also respect this floor. Equality forces
$q_{1i}=\eta$ for all three $i$. Since $J-I$ is invertible, with eigenvalues
$2,-1,-1$, it then forces $q_0=t(1,1,1)$. Thus the optimal positive-floor
face is a singleton and \eqref{eq:x} gives the full leximin allocation $x$.
No symmetry restriction on the feasible core was imposed.

Every level-$\ell$ pair changes from $z$ to $x$ by
$(2^\ell t,-2^\ell t)$, including the offset root pair. Each copied pair
changes by $(2^Lt,-2^Lt)$. Therefore
\[
 \|x-z\|_1=6t\left(\sum_{\ell=0}^{L}2^\ell+R2^L\right)
 =\frac{6((R+2)2^L-1)}{2^L+2}.
\]
Dividing by $T$ proves \eqref{eq:rho}. With $R=L^2$, both the amplified
payoff fraction and the share of copied pairs tend to one, giving $\rho\to2$.
Conversely, nonnegative vectors of total $T$ always satisfy
$\|x-z\|_1\le2T$, proving the sharp supremum.
For example, $L=4,R=16$ gives $\rho=287/198>4/3$.
Finally, every top $V$ has core maximum $1$, attained at $z$, but payoff
$2/(2^L+2)$ at $x$. This proves the individual conclusion separately.

\begin{proposition}\label{prop:dimension}
Let $S_n$ be the supremum for total player count $n$, allowing arbitrary sector
sizes. For $n\ge30$, put $k=\lfloor(n-12)/9\rfloor$. Then
\[
 0\le2-S_n\le
 \frac{2\lceil\log_2 k\rceil+4}{k+2}+\frac{2}{(k+2)^2}
 =O(\log n/n).
\]
\end{proposition}
\begin{proof}
Choose $L=\lceil\log_2 k\rceil$ and $R=k-L\ge0$, and add the remaining
at most eight players with zero incident values. Matching efficiency forces
these players to zero in the entire core, preserving both full selectors.
Subtracting \eqref{eq:rho} from two gives
\[
 2-\rho=\frac{2L2^L}{(2^L+2)(k+2)}
       +\frac4{2^L+2}+\frac2{(2^L+2)(k+2)}.
\]
Use $2^L\ge k$ and bound the first fraction by $2L/(k+2)$.
\end{proof}



\section{Additive triples and a separate individual-failure control}
\begin{proposition}\label{prop:additive}
For a complete three-sided market with $m$ players per sector and
additive nonnegative values $w_{ijk}=a_i+b_j+c_k$, full leximin satisfies
$x_i\ge H_i/3$ and $\|x-y\|_1/T\le4/3$ for every core point $y$.
Both constants are sharp across this class.
\end{proposition}
\begin{proof}
Subtract each sector's minimum from its coefficients. Write the resulting
nonnegative offsets as $b_{si}$, and let $S$ be the sum of the three original
minima. Complete-array nonnegativity implies $S\ge0$. Every perfect matching
has the same value, and each triple lies in one, so every triple inequality
is tight throughout the core. Subtracting these equalities shows that the
entire core is exactly
\[
 p_{si}=b_{si}+t_s,\qquad t_s\ge0,\qquad\sum_{s=1}^3t_s=S.
\]
Here $T=\sum_{s,i}b_{si}+mS$. This additive-core simplex is established
structure \cite{assortative}; its quantitative consequences are immediate. Each sector contains an offset-zero
player, so leximin uniquely sets all $t_s=S/3$ (also trivially when $S=0$).
Then $H_{si}=b_{si}+S\le3x_{si}$ and
\[
 \|x-y\|_1=m\sum_s|t_s-S/3|\le4mS/3\le4T/3.
\]
An offset-zero player makes the individual fraction sharp when $S>0$.
For $m\ge2$, the family $w_{ijk}=1+\mathbf1_{i=1}+\mathbf1_{j=1}$ has $T=m+2$,
leximin shifts $(1/3,1/3,1/3)$ and minimax shifts $(0,0,1)$. Its distance is
$4m/[3(m+2)]\to4/3$.
\end{proof}

Individual loss does not imply large selector disagreement. In the triple
recurrence construction, remove copies and root offsets, set every pair total
to one, remove root lower-bound triples, and change the level-one reverse
weight from three to one. The core is then $U_\ell=(J-I)^\ell U_0$,
$V_\ell=1-U_\ell$ with all dummy payoffs zero. A common floor $\eta$ implies
$\sum U_0\ge3\eta$ and $2^L\sum U_0\le3(1-\eta)$, so
$\eta\le1/(2^L+1)$. A ceiling $M$ implies $\sum U_0\ge3(1-M)$ and
$2^L\sum U_0\le3M$, so $M\ge2^L/(2^L+1)$. Both optima are unique at
$U_{\ell i}=2^\ell/(2^L+1)$. Thus $x=z$ although a top $V$ has
$x_i/H_i=1/(2^L+1)$.

\section{Separation from previously studied subclasses}
The recurrence array has genuine three-way interaction. Take sector coordinates
$A=(G_0,U_{00})$, $B=(G_1,V_{00})$, $C=(G_2,D_{00})$ and index the eight
corners by zeros and ones. Only $w_{111}=2$, $w_{110}=2$ and $w_{100}=1$
are nonzero. Its third mixed difference is
\[
 w_{111}-w_{110}-w_{101}-w_{011}+w_{100}+w_{010}+w_{001}-w_{000}=1.
\]
Every decomposition $w_{abc}=f_{ab}+g_{ac}+h_{bc}$ gives zero by cancellation.
Thus the family is outside even the pair-additive class, irrespective of
relabeling; it is not a pair market represented using decorative dummy agents.

A recent result for supermodular multisided arrays satisfying local monotonicity
 gives simultaneous attainment of all coordinate maxima within any one sector
\cite[Theorem~3.2]{extrema}. Our core violates this property for $L\ge2$.
Since $(J-I)^2=J+I$, every entry of $(J-I)^L$ is strictly positive. A top $V$
attains its maximum one only at $q_0=0$, while a root $U$ in the same sector
has maximum greater than one, as witnessed by a small positive feasible $q_0$.
Every sector contains both types. Hence no ordering of this array satisfies
that theorem's hypotheses. These are explicit exclusions of inspected
subclasses, not a proof of priority against all possible representations.

\section{Priority and interpretation}\label{sec:priority}
\paragraph{What is already established.}
Core/optimal-dual descriptions, matched-team complementarity, and the
multisided model are established background \cite{quint,atay}.
Sorted egalitarian selection is not new \cite{klijn,vazirani}.
Feasible pairwise equal splits are Lorenz-best, and exact assignment games
have coincident selections \cite{llerena,solymosi-raghavan}. These facts are
not candidate contributions here. Both extremal families are nonexact TU games:
a distinguished player has core minimum one despite singleton worth zero.

\paragraph{A direct dependence discovered in this audit.}
For the bipartite half-guarantee, start Vazirani's Algorithm~1 from an arbitrary
core point $y$ in its rational-weight domain. Its clock starts at the minimum
initial essential payoff $\Omega_0\ge0$; nonfixed payoffs move at
speed at most one, and upon permanent freezing at clock value $t_i$ the final
payoff obeys $x_i\ge t_i$. The update, freezing and floor invariants imply
\[
 |x_i-y_i|\le t_i-\Omega_0\le x_i-\Omega_0,
 \qquad 2x_i\ge y_i+\Omega_0\ge y_i.
\]
Initially fixed coordinates are unchanged. Relevant anchors in
\cite{vazirani} are Definition~19, Algorithm~1, Remark~1, Lemmas~5 and~9,
and Theorem~8 (printed pp.~17--18, 20, 22, 25, 27, 29--30). The inequality is our deduction from these statements, not a
located explicit theorem in that source. Accordingly the bipartite half
bound is best positioned as a short corollary with an elementary proof,
not as a wholly independent fairness discovery.

\paragraph{Candidate contribution and limits.}
The sharp finite distance bound, its attainment, the sparse bounded-integer
three-sided recurrence, and the combined comparison are the main candidates.
No equivalent quantitative theorem was located in the dated source audit.
That statement does not certify priority: original Quint and Lucas full texts
remain unavailable, and the preserved audit records exact access and citation
chain coverage. Existing assortative and local-monotonicity results impose
extra hypotheses; a three-way transportation-polytope representation also
would not automatically preserve player coordinates and full selectors.

The economic content is a worst-case dependence on nonadditive three-way
cooperation. The entire core has only three degrees of freedom.
The construction uses forced-zero dummies, exact constraints and
replicated top pairs. It does not show that these extremes are common in
observed markets or robust to arbitrary perturbations. The additive result
isolates a simple protected subclass; it makes no assertion about every
assortative market. The intended contribution is this sharp worst-case comparison, rather than
an empirical claim about matching markets. Independent specialist review of
correctness and priority is pending.

\paragraph{Verification and provenance.}
The accompanying portable package regenerates exhaustive bounded
$2\times2$ and $2\times2\times2$ calibrations, sampled
$3\times3\times3$ controls, and general-graph and additive checks.
It also includes independent sorted-rank improvement queries and a separate
raw-triple construction through 192 players. Run
\texttt{python verification/run\_checks.py}; the README records dependencies,
coverage and limitations. All-parameter conclusions rest on the proofs above,
not on finite enumeration. The claim/source ledger distinguishes established
results, our deductions, and unresolved priority questions. Codex assisted
with research development, drafting and computational verification; the
provenance file records source versions and package changes.
External specialist correctness and priority review is pending.

\begin{thebibliography}{9}
\small
\bibitem{vazirani} V.~V.~Vazirani. \emph{Fair Core Imputations for the Assignment
Game: New Solution Concepts and Efficient Algorithms}. arXiv:2402.11437v6,
20 November 2025 (proof version inspected). Listed for ACM EC 2026 on the
author's publication page, checked 5 September 2026.
\url{https://arxiv.org/abs/2402.11437};
\url{https://ics.uci.edu/~vazirani/}.
\bibitem{klijn} F.~Klijn et al. \emph{Neighbour Games and the Leximax Solution}.
Mathematical Methods of Operations Research 58 (2003), 191--208.
Linked text is the 1999 working-paper version.
\url{https://pure.uvt.nl/ws/files/533861/110.pdf}.
\bibitem{llerena} F.~Llerena. \emph{The pairwise egalitarian solution for the
assignment game}. Operations Research Letters 40 (2012), 84--88.
\url{https://doi.org/10.1016/j.orl.2011.11.011}.
\bibitem{solymosi-raghavan} T.~Solymosi and T.~E.~S.~Raghavan.
\emph{Assignment games with stable core}. International Journal of Game
Theory 30 (2001), 177--185. \url{https://doi.org/10.1007/s001820100072}.
\bibitem{quint} T.~Quint. \emph{The core of an $m$-sided assignment game}.
Games and Economic Behavior 3 (1991), 487--503.
\url{https://doi.org/10.1016/0899-8256(91)90017-9}.
\bibitem{atay} A.~Atay and M.~N\'u\~nez. \emph{A note on the relationship
between the core and stable sets in three-sided markets}.
Mathematical Social Sciences 98 (2019), 10--14.
\url{https://doi.org/10.1016/j.mathsocsci.2018.12.002}.
\bibitem{assortative} J.~Mart\'inez-de-Alb\'eniz, C.~Rafels and N.~Ybern.
\emph{Assortative multisided assignment games: The extreme core points}.
Games and Economic Behavior 120 (2020), 144--153.
\url{https://doi.org/10.1016/j.geb.2020.01.001}.
\bibitem{extrema} F.~J.~Mart\'inez-de-Alb\'eniz, C.~Rafels and N.~Ybern.
\emph{Maximum and minimum core payoffs for supermodular multisided assignment
markets with local monotonicity}. Discrete Applied Mathematics 379 (2026),
605--612. \url{https://doi.org/10.1016/j.dam.2025.10.006}.
\end{thebibliography}
\end{document}
